% GENERATED FILE. Edit canonical lessons, then run npm run generate:books.
% canonical-source-hash: fnv1a64:b76c71606e7c1944
% canonical-lessons: ES-C25-las-estaciones

\chapter{The Seasons}
\label{ch:spanish-seasons}

This chapter is generated from the canonical micro-lessons used by Language Ladder.
Each section stays independently resumable and preserves the authored prerequisite order.

\section[la primavera, el verano, el otoño, el invierno]{la primavera, el verano, el otoño, el invierno — three plain, one surprising}

\label{lesson:ES-C25-las-estaciones}



\begin{quote}
\textbf{Warm-up.} {[}PAUSE 2s{]} Three of Spanish's four seasons are exactly what you'd expect from Latin. The fourth hides a genuine surprise: its \textbf{meaning itself} drifted from one season to another.
\end{quote}

\begin{cousinweb}
\begin{itemize}
\raggedright
  \item \textbf{la primavera} ("spring") $\leftarrow$ Latin \emph{\textbf{prīma vēra}}, literally "\textbf{first spring}" (\emph{vēr}, "spring," + \emph{prīma}, "first"). Italian still says the identical phrase, \emph{primavera}, unchanged. (Old French once said \emph{primevère} too, before switching to \emph{printemps}.)
  \item \textbf{el otoño} ("autumn") $\leftarrow$ Latin \emph{\textbf{autumnus}} — the same root behind French \emph{automne} and English \emph{autumn} itself.
  \item \textbf{el invierno} ("winter") $\leftarrow$ Latin \emph{\textbf{hībernum}} (\emph{tempus}), "the \textbf{wintry} {[}season{]}" — the same root as French \emph{hiver} and English \textbf{hibernate}, "to spend the winter."
\end{itemize}
\end{cousinweb}

\begin{culture}
You might expect \emph{verano} ("summer") to come from Latin \textbf{aestas}, "heat" — the way French \emph{été} does. It doesn't. The standard account (per Spanish etymological dictionaries) traces \textbf{verano} to Latin \emph{\textbf{vērānum}} (\emph{tempus}), "\textbf{of spring}" — built on the very same \emph{vēr}, "spring," behind \emph{primavera}.

So how did a word meaning "spring-ish" end up naming \textbf{summer}? The honest answer is a \textbf{shift in meaning over time}, not a sound change — though the exact story is murkier than a tidy one-line explanation. The traditional, simplified account says that as \emph{primavera} ("first spring") rose to become the standard word for spring, \emph{verano}'s meaning drifted forward into the season right after, summer; the fuller picture likely also involves Vulgar Latin agricultural usage (a "vernal" grazing or working period that stretched from late spring into early summer), so the drift may not be purely a case of one word simply vacating a job for the other. Either way, two Spanish words both rooted in \emph{vēr}, "spring," ended up naming two different seasons.
\end{culture}

\subsection*{Guided Practice}
{[}PAUSE 1s{]}

\begin{itemize}
\raggedright
  \item {[}YOU SAY: "primavera, verano, otoño, invierno" — the four seasons{]}
  \item {[}YOU SAY: "primavera = prima vera, first spring"{]}
  \item {[}YOU SAY: the surprise — "verano ALSO comes from vēr, 'spring' — but its meaning moved to summer"{]}
  \item {[}YOU SAY: "invierno \textasciitilde{} hibernate"{]}
\end{itemize}

\begin{tcolorbox}[breakable,colback=teal!4,colframe=teal!35!black,title={Wrap-up recall}]
{[}PAUSE 3s{]} What does \textbf{primavera} literally mean? (\textbf{"First spring"} — \emph{prīma vēra}.) Does \textbf{verano} come from Latin \emph{aestas}, like French \emph{été}? (\textbf{No} — from \emph{vērānum}, "of spring," the same \emph{vēr} root as \emph{primavera}; its meaning \textbf{shifted} toward summer over time — traditionally linked to \emph{primavera} taking over the job of naming spring, though the fuller history is murkier than that alone.) What English word is \textbf{invierno} the cousin of? (\textbf{Hibernate}.)
\end{tcolorbox}
